\section{技术专家二：世界模型、强化反馈与交通闭环}

上一章讲 VLA 是什么，本章解释它如何训练和验证。李想提到，理想会使用世界模型生成数据来训练，让模型从 A 点开到 B 点，并用舒适性、碰撞、交通规则等反馈判断结果。这里可以看到一个很明确的模式：真实车队提供数据，世界模型生成和考试，VLA 在数据和反馈中更新，最后再部署回车辆。

\begin{figure}[H]
\centering
\includegraphics[width=0.96\textwidth]{world-model-loop.png}
\caption{交通世界模型闭环：真实驾驶数据、仿真生成和 VLA 训练互相推动。自制概念图，依据 00:47:30--00:48:10 与 01:10:06--01:10:48 对谈内容整理。}
\end{figure}

\begin{knowledgebox}{读图：世界模型有三种角色}
第一，它可以当考试系统，测试 VLA 在不同交通场景中的表现。第二，它可以生成训练数据，覆盖真实道路中难以高频采集的场景。第三，它未来可能成为无人车运营系统的一部分，帮助没有司机的车辆理解和预测交通世界。
\end{knowledgebox}

\subsection{世界模型不是“画面生成器”}

在自动驾驶语境里，世界模型不只是生成漂亮视频。它必须对动作后果敏感：如果车辆变道、加速、避让或等待，交通世界会怎样演化。它还要对安全和规则敏感：碰撞、急刹、闯灯、压线、插队、舒适性都不只是视觉问题，而是任务结果问题。

\begin{importantbox}{世界模型的最小公式}
可以把交通世界模型粗略写成：
\[
s_{t+1} = f(s_t, a_t, c_t).
\]
其中，\(s_t\) 表示当前交通状态，\(a_t\) 表示车辆动作，\(c_t\) 表示场景约束和规则，\(s_{t+1}\) 表示下一时刻状态。真正难的是让 \(f\) 对真实道路的长尾场景足够可靠。
\end{importantbox}

\subsection{验证成本与模型规模}

李想提到验证成本从每一万公里 18 万降到 4000 元，这说明 VLA/世界模型路线不只是模型能力问题，也会改变验证经济学。自动驾驶如果只能靠真实道路堆里程，成本会非常高；世界模型和仿真可以降低某些测试成本，但最终仍需要真实世界验证兜底。

\begin{warningbox}{仿真不能直接替代真实世界}
世界模型可以降低探索和验证成本，但不能把真实世界完全替换掉。仿真中没有覆盖的长尾行为、传感器噪声、道路施工、极端天气和人的非理性动作，仍然可能在真实道路中出现。
\end{warningbox}

\subsection{L3、L4 与车端算力上限}

前面讨论了世界模型如何降低训练和验证成本，本节把问题落到自动驾驶等级与端侧约束。访谈中，李想把 L3/L4 能力和云端、端侧模型规模联系起来。他认为当前一代算力大致对应 L3 能力，L4 还需要更多能力和更多条件。这个判断的教育意义在于：自动驾驶等级不是单纯软件开关，而是模型规模、端侧算力、验证体系、法规和运营责任共同决定。

\begin{knowledgebox}{术语消化：L3 与 L4}
L3 通常意味着在特定条件下系统可以承担驾驶任务，但需要人类在系统请求时接管；L4 则意味着系统可在限定运营设计域内不依赖人类接管。二者的差别不仅是模型分数，更是责任边界、冗余设计、法规和运营系统的差别。
\end{knowledgebox}

\subsection{本章小结}

世界模型把 VLA 从“会开车的模型”推进为“可训练、可考试、可运营的系统”。它降低验证成本，但不能消除真实世界的不确定性；它提高上限，但也要求更强的算力、数据和安全体系。

